\section{引言}\label{sec:intro}

有限维空间上的线性变换, 其谱即特征值之全体, 至多含 $n$ 个点. 无穷维空间上的有界算子则不然: 谱可以是连续统, 亦可以不含任何特征值. 紧算子介于二者之间, 其谱的结构与矩阵十分接近, 这正是积分方程理论得以建立的基础 \cite{riesz1918,fredholm1903}.

% 下一段是 AI 润色演示所用的原稿, 措辞刻意保留了一些冗余.
在这篇文章里面, 我们主要是想要去讨论一下紧算子的谱的一些性质, 并且我们还会通过一个具体的积分算子的例子, 来对 Nyström 方法计算出来的特征值的精度进行一个比较详细的考察和分析.

\Cref{sec:prelim} 介绍紧算子的定义与基本性质; \cref{sec:spectral} 证明 Riesz--Schauder 定理及紧自伴算子的谱分解定理; \cref{sec:numerics} 给出数值实验. 记号与证明细节见 \cref{app:proofs}.
